% !PRINTED 370
\weberpage{370}
% Source-faithful three-stroke Cayley configuration sign.
\providecommand{\weberthreeup}{\mathord{\tikz[baseline=-0.2ex,x=0.23em,y=0.23em,line width=0.45pt]{\draw (0,0)--(-1,1) (0,0)--(0,1.08) (0,0)--(1,1);}}}
%<<CONT
contains, and which forms an azygetic pair with (1). It must
therefore be of the form:
\weq{2}{x_1 y_2,\ x_3 z_3,\ x_4 z_4,\ x_5 z_5,\ x_6 z_6,\ x_7 z_7,}
and accordingly $y_2$ is also to be denoted by $[2\,1]$, q. e. d.

The $z$ in (2) are different from the $y$ in (1), and since, e.g.,
$z_3 = [2\,3]$, it follows that in general two $[\mu\,\nu]$ which do not
agree in both figures $\mu, \nu$ are different from one another.
Since twenty-one pairs can be formed from seven figures, one
obtains in this way all the double tangents, and each only once.

From this mode of representation it also follows that no
azygetic systems of more than seven double tangents exist;
for if we add to the seven $x$ any further double tangent, for
which, the figures 1 to 7 being equally entitled, we may choose
say $y_2 = [1\,2]$, then $y_2 x_1 x_2$ is a syzygetic triple.

It is expedient to introduce an eighth figure 8, and to
denote the elements of the original system of seven not by the
single figures, but by the pairs
\[
[1\,8],\ [2\,8],\ [3\,8],\ [4\,8],\ [5\,8],\ [6\,8],\ [7\,8]
\]
where it shall then also hold that $[1\,8] = [8\,1]$, and so forth.
Then all 28 double tangents are denoted uniformly by the pairs
$[\mu\,\nu]$, in which $\mu$ and $\nu$ denote two different figures of the
sequence 1 to 8.

Here it is very helpful for clarity to employ an intuitive
notation\footnote{After \gesperrt{Cayley}; cf. \gesperrt{Salmon}, ``Higher plane curves''.}. One indicates
a double tangent $[\mu\,\nu]$ by a single stroke $\mid$, at
whose ends one imagines the two figures $\mu, \nu$ placed.
Then two strokes without a common point, $\mid\mid$, denote two
double tangents in whose designation $[\mu\,\nu]$ no common
figure occurs, and two strokes issuing from one point, $\vee$, two
double tangents which have a common figure in their symbols $[\mu\,\nu]$.
After this the more complicated signs that we employ later are
self-explanatory. For a triple of double tangents we have, e.g., the
following five signs $\mid\mid\mid$, $\sqcap$, $\triangle$, $\weberthreeup$, $\vee\mid$.
